\begin{abstract}
Shy and Stenbacka (2007) study regulation of banks that offer risky and liquid deposit accounts. I re-examine the reduced bank-profit problem used to derive their deposit-rate and fee formulas. Direct differentiation of the stated reduced objective yields additional terms and a different Hessian, so the displayed best response does not solve that objective. I derive the corrected reduced best response, symmetric rate and fee, and comparative statics. In particular, both prices generally depend on the banks' investment return. I also identify a separate constrained-optimization issue in the regulator's problem: when the liquidity value exceeds the investment return, optimization along the binding reserve boundary selects full risky-account coverage with sufficient reserves rather than the reported interior mix. The aggregate-welfare identity itself survives. The analysis is deliberately limited to the source's reduced problem and symmetric implementation; it does not complete unspecified off-path states of the original two-instrument game.
\end{abstract}

\noindent\textbf{Keywords:} bank regulation; liquidity provision; deposit competition; reserve requirements; mathematical correction.\\
\textbf{JEL classification:} G21; G28; L13; C72.

\section{Introduction}
\citet{ShyStenbacka2007} develop a differentiated-bank model in which a regulator chooses both a reserve requirement and the maximum fraction of risky deposit accounts. Depositors differ in the probability of needing early liquidity. Banks compete for depositors using an interest rate on risky accounts and a fee on liquid accounts. The model's central policy mechanism is that regulation can induce self-selection across account types and thereby change the cost of providing liquidity. The paper belongs to a broader line of work by the same authors on narrow banking, deposit-market competition, and liquidity regulation \citep{ShyStenbacka2000,ShyStenbacka2008,ShyStenbacka2019,ShyStenbacka2024}.

This note reassesses two mathematical steps in that model. The first concerns the reduced bank-profit objective used to derive the competitive risky-account rate. In the complete author manuscript corresponding to the 2007 article, Eq.~(7) is a one-dimensional reduced optimization problem obtained after imposing the marginal depositor's account-indifference condition. Direct differentiation of that objective produces terms that are absent from the displayed Eq.~(8). The second derivative is also different from the expression reported immediately below Eq.~(7). The corrected objective is strictly concave, so its maximizer is simple to characterize, but the resulting rate and fee differ from Eqs.~(11)--(12). Most notably, the corrected symmetric prices generally depend on the banks' investment return $r$, contrary to Proposition~1(c).

The second issue is logically separate. The aggregate-welfare identity in Eq.~(14) and the two partial derivatives in Eq.~(15) can be reproduced. However, when $v>r$, the optimal reserve condition places the planner on the moving boundary $\rho=\delta/2$. Optimizing $\delta$ along that boundary requires a total derivative. Doing so yields $\mathrm dW/\mathrm d\delta=r(1-\delta)$, so for $v>r>0$ the source candidate is dominated by $(\delta,\rho)=(1,1/2)$. The exact welfare difference is $r(v-r)^2/[2(r+v)^2]$. Thus the welfare accounting survives, while the reported mixed-banking optimizer in Proposition~2(a) does not.

The scope is intentionally narrow. I solve completely the reduced one-dimensional problem actually used to obtain the published rate and fee formulas, and I solve the regulator's constrained welfare problem including reserve clipping and the endpoints $\delta=0$ and $\delta=1$. I do not claim a complete global equilibrium characterization of asymmetric off-path states. The original model does not fully specify how market shares, reserve feedback, or account-type participation should continue after sufficiently large unilateral deviations. Strict concavity of the reduced objective therefore should not be read as uniqueness of the unrestricted original banking game.

\paragraph{Source-version note.}
The equation-level calculations in this note are keyed to the complete author/institution-hosted manuscript corresponding to \citet{ShyStenbacka2007}. Bibliographic metadata for the Wiley Version of Record (VOR) are independently confirmed, but I have not obtained equation-level access to the full VOR. Accordingly, references below to Eqs.~(4)--(17) and Propositions~1--2 refer to that complete manuscript, and any statement about the journal VOR is conditional on the corresponding equations being unchanged there. This qualification can be removed only after direct VOR comparison.
